\begin{abstract}
  This essay presents the theoretical foundation of a four-part research project
  developing a symbolic and structural theory of human encounter, based on the concept of \textit{life rays}.
  Each person is understood as an individual trajectory in space and time, whose development is shaped by points of contact with others.
  These intersections generate \textit{fragments} -- units of memory or imprint -- which remain encoded along one's biographical timeline.

  The theory seeks to explain how emotional depth, social proximity, or systemic connectedness between individuals emerges -- or fails to emerge.
  It links phenomenological thinking (Merleau-Ponty), narrative identity (Ricoeur), systems theory (Luhmann), and informatic structure (network theory)
  into an interdisciplinary conceptual model.\footnote{For contextual orientation only: parallels can be drawn to Maurice Merleau-Ponty’s phenomenology of perception,
  Paul Ricoeur’s concept of narrative identity, and Niklas Luhmann’s systems theory, as well as to approaches in network theory and distributed systems in computer science.
  The present model, however, was developed independently and not derived from any of these frameworks.}

  As the first part of the project, this paper focuses on the conceptual and philosophical architecture of the model.
  Subsequent parts will address its formal, empirical, and computational realizations - ranging from mathematical characterization and data modeling to applied simulations and AI design.

  The aim is not to reduce interpersonal complexity, but to render it graspable in a structured and transferable way.
  The model is intended to be applicable in psychological, social, and technological contexts - and it is presented as an open invitation for further development by other readers, thinkers, system designers, and builders.
\end{abstract}

\vspace{1em}
\begin{center}
  \textbf{Keywords:}
  life rays, shared space, own space, fragment chain, resonant asymmetry,
  humor theory, cognitive incongruity, intersubjectivity, relational dynamics,
  encounter, memory density
\end{center}
